\documentclass[10pt]{beamer}

\usetheme[progressbar=frametitle]{metropolis}
\usepackage{booktabs}
\usepackage{amsmath}
\usepackage{graphicx}
\usepackage{tikz}

\title{Improving VLA Models with\\Contrastive Representation Learning}
\subtitle{DL2 Project Update — Week 2, May 13 2026}
\author{Simon Bruno · Finley Helms · Sem IJsendijk · Lev Hulsbergen · Mitch Boontjes · Wes André}
\institute{University of Amsterdam}
\date{}

\begin{document}

\maketitle

% ── SLIDE 1: What we did ────────────────────────────────────
\begin{frame}{What We Did This Week}

\textbf{Switched to LIBERO benchmark} (from RoboCasa)
\begin{itemize}
  \item 4 suites: Spatial · Object · Goal · Long (10 tasks each)
  \item 2 cameras (exterior + wrist) — ViewCutoff works correctly
  \item Faster iteration: 12h training vs 48h on RoboCasa
\end{itemize}

\vspace{0.8em}

\textbf{RS-CL reproduction} (60K steps, frozen VLM, all 4 suites)
\begin{itemize}
  \item Baseline and RS-CL trained with paper-matching setup
  \item Fixed: token alignment, $\beta=1.0$, 256 tokens/view
\end{itemize}

\vspace{0.8em}

\textbf{Modality ablations} (10K steps, Long only)
\begin{itemize}
  \item Depth maps extracted (547K PNGs, Depth-Anything-V2)
  \item Tested: depth, action, EE velocity as soft-label distances
\end{itemize}

\end{frame}

% ── SLIDE 2: Main results ───────────────────────────────────
\begin{frame}{RS-CL Reproduction Results}

\begin{table}
  \centering
  \small
  \begin{tabular}{lccccc}
    \toprule
    \textbf{Method} & \textbf{Spatial} & \textbf{Object} & \textbf{Goal} & \textbf{Long} & \textbf{Avg.} \\
    \midrule
    Paper baseline   & 98.2 & 99.4 & 97.2 & 87.8 & 95.7 \\
    Paper + RS-CL    & 98.4 & 98.6 & 98.2 & \textbf{90.4} & 96.4 \\
    \midrule
    Our baseline     & 97.6 & 93.2 & 93.6 & 85.8 & 92.6 \\
    Our + RS-CL      & 97.0 & 93.2 & 91.8 & \textbf{87.6} & 92.4 \\
    \bottomrule
  \end{tabular}
\end{table}

\vspace{0.5em}

\begin{itemize}
  \item RS-CL improves \alert{Long tasks by +1.8\%} (paper: +2.6\%)
  \item Spatial/Object near ceiling — no room for improvement
  \item Matches the paper's pattern: gains on harder, multi-step tasks
\end{itemize}

\end{frame}

% ── SLIDE 3: Modality ablation ──────────────────────────────
\begin{frame}{Modality Ablation — Depth Works!}

\begin{table}
  \centering
  \begin{tabular}{lcccc}
    \toprule
    \textbf{Soft-label distance} & $w_\text{depth}$ & $w_\text{action}$ & $w_\text{vel}$ & \textbf{Long (\%)} \\
    \midrule
    Proprio only (control) & 0 & 0 & 0 & 90.2 \\
    \alert{+ Depth}        & 1 & 0 & 0 & \alert{\textbf{92.0}} \\
    + Action               & 0 & 1 & 0 & eval pending \\
    + Velocity             & 0 & 0 & 1 & pending \\
    \bottomrule
  \end{tabular}
\end{table}

\vspace{0.5em}

\begin{itemize}
  \item \textbf{Depth: +1.8\%} over proprio-only RS-CL on Long
  \item Depth adds 3D geometry the 2D vision encoder can't see
  \item Action: no change in training proxy (hpc); eval pending
  \item Velocity \& 30K depth evals running overnight
\end{itemize}

\end{frame}

% ── SLIDE 4: Key insight ────────────────────────────────────
\begin{frame}{Key Insight: Contrastive Loss Dynamics}

\textbf{Observation:} $\mathcal{L}_\text{CL}$ stays flat at $\sim\!\log B$ throughout training

\vspace{0.5em}

\begin{itemize}
  \item $\cos(z, \hat{z})$ collapses to 0.99 by step 3K
  \item Self-attention adapter reconstructs masked camera view
  \item \textbf{But RS-CL still helps} — gradients reshape adapter representations (h) to align with proprio
  \item We added \textbf{hpc} (h-proprio correlation) metric to track this
\end{itemize}

\vspace{0.8em}

\textbf{Implication:} RS-CL works through gradient-level representation shaping, not explicit contrastive loss minimization. The paper doesn't report $\mathcal{L}_\text{CL}$ trajectories.

\end{frame}

% ── SLIDE 5: Next steps ─────────────────────────────────────
\begin{frame}{Next Steps}

\textbf{This week (May 13--20):}
\begin{itemize}
  \item Finish depth/velocity eval results
  \item Full 60K depth run on all 4 suites (if 10K results hold)
  \item Implement GRAM as alternative contrastive objective
  \item Start writing discussion + analysis sections
\end{itemize}

\vspace{0.8em}

\textbf{Week 3 (May 20--27):}
\begin{itemize}
  \item Final results table with all ablations
  \item Finalize 4-page report
  \item Clean repo + reproducibility check
\end{itemize}

\vspace{0.8em}

\textbf{Open question:} Should we also run on RoboCasa for cross-benchmark validation, or focus on LIBERO depth?

\end{frame}

\end{document}
